\documentclass[10pt,a4paper]{amsart}
\usepackage[margin=19mm]{geometry}
\usepackage{amsmath,amssymb,mathtools}
\usepackage[scr=rsfs,cal=euler]{mathalfa}
\usepackage{xfrac}
\usepackage[unicode,hidelinks,bookmarks=false]{hyperref}
\newcommand{\ie}{i.e.\ }
\newcommand{\cf}{cf.\ }
\newcommand{\resp}{respectively\ }
\newcommand{\Subsection}{\subsection}
\providecommand{\nobreakdash}{\nobreak\hbox{-}}
\setlength{\emergencystretch}{2em}
\multlinegap0pt
\usepackage{amssymb}
\usepackage{tikz}
\newtheorem{lemma}{Lemma}
\newtheorem{theorem}{Theorem}
\title{Betti Numbers of Sequentially Cohen-Macaulay Co-Chordal Graphs and Their Applications}
\author{Mohammed Rafiq Namiq}
\date{Source: arXiv:2606.22513v2 (CC BY 4.0)}
\begin{document}
\maketitle

\noindent\emph{Source and licence.} Betti Numbers of Sequentially Cohen-Macaulay Co-Chordal Graphs and Their Applications. Version arXiv:2606.22513v2 (CC BY 4.0), distributed by arXiv under the Creative Commons Attribution 4.0 International licence, http://creativecommons.org/licenses/by/4.0/, as recorded in the arXiv licence metadata for this exact version and shown on the abstract page https://arxiv.org/abs/2606.22513v2. Full licence notice: \texttt{licenses/arxiv-2606.22513-CC-BY-4.0.txt}.\par\smallskip

\noindent\textit{Editorial scope.} Counted unit: the article's theorem thm:master\_nilpotent (source0.tex lines 673--712). The article's complete original proof of it is delivered with it (source0.tex lines 713--849, one \texttt{proof} environment of 137 lines). No mathematics is written, repaired or re-derived: the statement and the proof are the article's own lines, in the article's own notation, and no retained line is substituted or re-flowed.  Retained uncounted as the source-local context the proof needs: lines 90--92; lines 318--320; lines 331--345.  Nothing was omitted from any retained span, so the package carries no omission record.  No retained line contains a citation, so the wrapper carries no bibliography and no bibliography entry.  No retained line contains a figure environment, an image-inclusion command or drawing code, so no image is delivered and the package declares no asset dependency.  Editorial formatting only, in addition: an AMS \texttt{amsart} preamble that restates the article's own declarations and macros used by the retained lines, namely the environments \texttt{lemma}, \texttt{theorem} on separate counters, together with the standard packages those lines need.  The retained environments print in this document as Lemma 1, Theorem 1 in order of appearance, whereas the article prints the same items with the numbering of its own full document.  The complete native source of the article as distributed in the arXiv e-print for this exact version is bundled unchanged as \texttt{sources/203350/source0.tex}.
	\begin{lemma}\label{lem:ridge-obstruction}
Let $G$ be co-chordal.  If $\operatorname{Ind}(G)$ has two facets, each of which shares no ridge with any other facet, then $S/I(G)$ is not sequentially Cohen--Macaulay.
\end{lemma}

	\begin{theorem}\label{thm:split_is_scm}
Every split graph is co-chordal and sequentially Cohen--Macaulay.
\end{theorem}

	\begin{theorem}\label{thm:split_betti}
Let $G$ be a split graph with $V(G)=C\sqcup U$, where $C=\{x_1,\ldots,x_n\}$ and $n\ge1$.  Put
\[
n_k=|N_G(x_k)\cap U|.
\]
Then
\[
\beta_i(S/I(G))
=
\sum_{k=1}^{n}\binom{n+n_k-1}{i}
-
\binom{n}{i+1}
\qquad(i\ge1).
\]
\end{theorem}

\begin{theorem}\label{thm:master_nilpotent}
Let $R\cong\prod_{i=1}^{k}R_i$ be a finite direct product of commutative Artinian chain rings.  Let $\mathfrak m_i$ be the maximal ideal of $R_i$, let $a_i$ be its nilpotency index, and put $q_i=|R_i/\mathfrak m_i|$.  Set
\[
r=\bigl|\{i:a_i\ge3\}\bigr|.
\]
Then the following conditions are equivalent.
\begin{enumerate}
\item $\Gamma_{\operatorname{Nil}}(R)$ is co-chordal and $S/I(\Gamma_{\operatorname{Nil}}(R))$ is sequentially Cohen--Macaulay.
\item $\Gamma_{\operatorname{Nil}}(R)$ is a threshold graph.
\item At most one factor $R_i$ has $a_i\ge3$, equivalently $r\le1$.
\end{enumerate}

Suppose that these conditions hold and that exactly one factor, say $R_1$, has $a_1\ge3$.  Put
\[
M=\prod_{i=2}^{k}q_i^{a_i-1}.
\]
Then
\begin{equation}\label{eq:nil_product}
\beta_i(S/I(\Gamma_{\operatorname{Nil}}(R)))
=
\sum_{t=\lceil a_1/2\rceil}^{a_1-1}|W_t|
\binom{|C|+n_t-1}{i}
-
\binom{|C|}{i+1},
\end{equation}
where
\[
|W_t|=
\begin{cases}
M(q_1^{a_1-t}-q_1^{a_1-t-1}),&t<a_1-1,\\
Mq_1-1,&t=a_1-1,
\end{cases}
\]
and
\[
|C|=Mq_1^{\lfloor a_1/2\rfloor}-1,
\qquad
n_t=M(q_1^t-q_1^{\lfloor a_1/2\rfloor}).
\]
\end{theorem}

\begin{proof}
For each $i$, choose a generator $\pi_i$ of $\mathfrak m_i$ and define, for $x_i\in\mathfrak m_i$,
\[
\nu_i(x_i)=
\begin{cases}
t,&x_i\in\mathfrak m_i^t\setminus\mathfrak m_i^{t+1},\\
a_i,&x_i=0.
\end{cases}
\]
Since $R_i$ is a chain ring,
\begin{equation}\label{eq:valuation-zero-product}
x_iy_i=0
\quad\Longleftrightarrow\quad
\nu_i(x_i)+\nu_i(y_i)\ge a_i.
\end{equation}
The vertex set of $G:=\Gamma_{\operatorname{Nil}}(R)$ is
\[
V(G)=\left(\mathfrak m_1\times\cdots\times\mathfrak m_k\right)\setminus\{0\}.
\]
Consequently, two distinct vertices $x$ and $y$ are adjacent in $H:=\overline G$ if and only if
\begin{equation}\label{eq:nil-complement-adjacency}
\nu_i(x_i)+\nu_i(y_i)<a_i
\quad\text{for at least one }i.
\end{equation}
A coordinate with $a_i\le2$ never contributes an edge to $H$.

Assume first that condition~(3) holds.  If $r=0$, the product of any two nilpotent elements is zero, so $G$ is complete and therefore threshold.  If $r=1$, relabel the factors so that $a_1\ge3$.  Equation~\eqref{eq:valuation-zero-product} gives
\[
xy=0
\quad\Longleftrightarrow\quad
\nu_1(x_1)+\nu_1(y_1)\ge a_1,
\]
because the remaining coordinates always have zero product.  Thus $G$ has a threshold-weight representation with vertex weight $\nu_1(x_1)$ and threshold $a_1$.  This proves $(3)\Rightarrow(2)$.

Threshold graphs are closed under complementation, and every threshold graph is split.  Hence $G$ is co-chordal, and Theorem~\ref{thm:split_is_scm} implies that $S/I(G)$ is sequentially Cohen--Macaulay.  Therefore $(2)\Rightarrow(1)$.

We prove $(1)\Rightarrow(3)$.  Assume that $H$ is chordal and that $S/I(G)$ is sequentially Cohen--Macaulay.  Suppose, toward a contradiction, that $r\ge2$, and set
\[
s=\bigl|\{i:a_i\ge4\}\bigr|.
\]
We first use chordality to restrict the possible active nilpotency indices.

If $s\ge2$, choose active coordinates with indices $a,b\ge4$, and put
\[
u=\left\lceil\frac a2\right\rceil,
\qquad
v=\left\lceil\frac b2\right\rceil.
\]
Choose four vertices whose valuations in these two coordinates are
\[
(1,b-1),\qquad (u,v),\qquad (a-1,1),\qquad (u,v),
\]
using distinct elements for the two vertices with valuation pair $(u,v)$ and setting all other coordinates equal to zero.  By~\eqref{eq:nil-complement-adjacency}, consecutive vertices are adjacent in $H$, while the opposite pairs are nonadjacent.  They therefore induce a $4$-cycle, contrary to chordality.

If $s=1$ and $r\ge3$, choose active coordinates with indices $a\ge4,3,3$ and vertices with valuation triples
\[
(1,2,2),\qquad
\left(\left\lceil\frac a2\right\rceil,1,2\right),\qquad
(a-1,1,1),\qquad
\left(\left\lceil\frac a2\right\rceil,2,1\right).
\]
Again, consecutive pairs and only consecutive pairs satisfy~\eqref{eq:nil-complement-adjacency}; hence these vertices induce a $4$-cycle.

If $s=0$ and $r\ge4$, choose four active coordinates, all of nilpotency index $3$.  For a vertex $x$, let its shallow support be the set of these coordinates in which $\nu_i(x_i)=1$, taking valuation $2$ in the remaining active coordinates.  Vertices with shallow supports
\[
\{1,2\},\qquad \{2,3\},\qquad \{3,4\},\qquad \{4,1\}
\]
induce a $4$-cycle, because two such vertices are adjacent in $H$ exactly when their shallow supports intersect.

It follows that, under the assumption $r\ge2$, chordality leaves only the following possibilities:
\begin{enumerate}
\item $r\in\{2,3\}$ and every active index equals $3$;
\item $r=2$, one active index is at least $4$, and the other equals $3$.
\end{enumerate}
We show that both possibilities contradict sequential Cohen--Macaulayness.

In the first case, for every active coordinate $i$, define
\[
L_i=\{x\in V(H):\nu_i(x_i)=1\}.
\]
The set $L_i$ is a clique in $H$.  It is maximal: if $z\notin L_i$, choose a vertex $e$ whose $i$th coordinate has valuation $1$ and whose other coordinates are zero.  Then $e\in L_i$ and $e$ is nonadjacent to $z$.  Let $E_i$ be the set of all such vertices $e$.  Since
\[
|\mathfrak m_i\setminus\mathfrak m_i^2|=q_i^2-q_i\ge2,
\]
the set $E_i$ contains at least two vertices.  Every neighbor of a vertex in $E_i$ belongs to $L_i$; therefore every maximal clique different from $L_i$ omits all of $E_i$.  Thus $L_i$ shares no ridge with any other maximal clique.  Since there are at least two active coordinates, the clique complex of $H$, equivalently $\operatorname{Ind}(G)$, has at least two ridge-isolated facets.  This contradicts Lemma~\ref{lem:ridge-obstruction}.

In the second case, relabel so that $a_1=a\ge4$ and $a_2=3$, and put
\[
L=\{x:\nu_2(x_2)=1\},
\qquad
f=\left\lfloor\frac{a-1}{2}\right\rfloor,
\qquad
h=\left\lfloor\frac a2\right\rfloor.
\]
The set $L$ is a maximal clique.  Indeed, if $z\notin L$, then every vertex whose first coordinate is zero and whose second coordinate has valuation $1$ is nonadjacent to $z$.  This block contains at least two vertices, and every maximal clique distinct from $L$ omits it.  Hence $L$ is ridge-isolated.

Define
\[
Q=
\{x\notin L:\nu_1(x_1)\le f\}
\ \cup\
\{x\in L:\nu_1(x_1)\le h\}.
\]
The set $Q$ is a clique.  Two vertices in the first part are adjacent through the first coordinate because $2f<a$; two vertices in the second part are adjacent through the second coordinate; and a vertex from each part is adjacent through the first coordinate because $f+h<a$.

The clique $Q$ is maximal.  If $z\in L\setminus Q$, then $z$ is nonadjacent to every vertex outside $L$ whose first-coordinate valuation is $f$.  If $z\notin L\cup Q$, then $z$ is nonadjacent to every vertex in $L$ whose first-coordinate valuation is $h$.  Both boundary blocks contain at least two vertices.  Moreover, every maximal clique different from $Q$ contains a vertex outside $Q$ and consequently omits one of these two entire blocks.  Hence $Q$ is ridge-isolated.  The two ridge-isolated facets $L$ and $Q$ contradict Lemma~\ref{lem:ridge-obstruction}.

Both remaining cases are impossible; therefore $r\le1$.  This proves $(1)\Rightarrow(3)$ and completes the equivalence.

It remains to prove the Betti formula.  Assume that $R_1$ is the unique factor with $a_1\ge3$.  Since $|\mathfrak m_i|=q_i^{a_i-1}$, the last $k-1$ coordinates contribute the multiplicity
\[
M=\prod_{i=2}^{k}q_i^{a_i-1}.
\]
The vertices with first-coordinate valuation at least $\lceil a_1/2\rceil$, together with the vertices whose first coordinate is zero, form a clique $C$; the remaining vertices form an independent set.  Their cross-neighborhoods are nested by valuation.

In $R_1$, the valuation-$t$ layer has cardinality
\[
q_1^{a_1-t}-q_1^{a_1-t-1}.
\]
For $t<a_1-1$, this gives
\[
|W_t|=M(q_1^{a_1-t}-q_1^{a_1-t-1}).
\]
At $t=a_1-1$, the $M-1$ vertices with first coordinate zero are universal and join the $M(q_1-1)$ ordinary layer vertices, so
\[
|W_{a_1-1}|=Mq_1-1.
\]
Summing the clique layers gives
\[
|C|=Mq_1^{\lfloor a_1/2\rfloor}-1.
\]
A clique vertex of valuation $t$ is adjacent to exactly
\[
n_t=M(q_1^t-q_1^{\lfloor a_1/2\rfloor})
\]
vertices in the independent set.  Applying Theorem~\ref{thm:split_betti} proves~\eqref{eq:nil_product}.
\end{proof}

\end{document}
